% ULSBS songbook template (A5)

% If you want different paper size, rename this file to not include _A5 and
% change the paper option to e.g. "paper={a4paper}".
%
% See ulsbs/WRITING_SONGBOOKS.md for document class options and song syntax.

\documentclass[
  variants={default}, % Add lyric, charango or bassclef to make them available.
  paper={a5paper},
  maintitle={My Songbook},
  author={Songbook Author},
  subject={Songbook}
]{ulsbs-songbook}

% Optional source cover icons (can have multiple: [foo][bar][etc]).
\ulSetVariableCoverIcons[ulsbs-tag-icon-chords_EN_512x512px]

\renewcommand{\ulImprintPageFootnote}{%
  This is a songbook created with ULSBS (Unilaiva Songbook System).
}

% For visibility switches and fonts, see the layout and style reference in
% ulsbs/WRITING_SONGBOOKS.md. Put preamble settings here, after the class loads.
% Select variants through the compiler configuration, not in this source.
% Use \ulLocalShow{feature}{true|false} for grouped requests in song content;
% a local true cannot override a book-wide false.

\begin{document}

% Cover page. Give PDF file as argument; if left empty, only songbook's name is
% displayed):
\ulCoverPage{}

% Optional: imprint page
\ulImprintPage
% Optional: table of contents
\ulToc


% Main chapter (you may have multiple of these)
\ulMainChapter{Main Chapter}{blue}

% All songs must be within "ulSongs"; the required list routes index records.
% Title, author and tag indices are generated; use {} for no automatic records.
\begin{ulSongs}{titleidx,authidx,tagidx}

% Example song
\beginsong{Example Song}[id={example-song},by={Example Author},key={C}]
  \beginverse
    |\[C]This is an example |song line.
    Another |\[G7]line of |\[C]lyrics.
  \endverse
\endsong

\end{ulSongs}


\end{document}
